\copyrightTim

\begin{frame}[fragile]{Sections}
Het commando \mintinline{tex}{\section{...}} maakt een heading (titel, kop, tussenkopje). Deze headings worden automatisch genummerd.
Hierbij typ je bij de ... dus je titel. 
\pause
\newline
Je kan met \mintinline{tex}{\subsection{...}} en \mintinline{tex}{\subsubsection{...}} onderkopjes maken.
\pause
\newline
Wil je geen nummering in je kopjes, gebruik dan \mintinline{tex}{\paragraph{}}.
\pause
\newline
voorbeeld
\begin{minted}[highlightlines={3,6}, fontsize=\small]{tex}
\documentclass[a4paper]{article}
\begin{document}
\section{pakkende titel}
Het bedenken van een titel kan erg ingewikkeld zijn,
het moet een samenvatting zijn van wat er in de tekst wordt gezegd.
\subsection{subtitel, ook pakkend}
Hiervoor kan het fijn zijn om subtitels te gebruiken, hiermee kan je 
aangeven waar specifieke stukken van de tekst over gaan
\end{document}
\end{minted}


probeer dit zelf uit in overleaf\par
\end{frame}